\chapter{Transición de nivel}

\section{Cambio de nivel}
\begin{definition}
Se llama transición de nivel al paso desde una formulación estructural ya legible hacia otra formulación que incrementa su capacidad de organización, comparación o tratamiento sin perder dependencia canónica respecto del núcleo basal integrado en Tomo I.
\end{definition}

\begin{observation}
La transición de nivel no equivale a ruptura arbitraria. Su función es ordenar aumentos de capacidad formal sin desanclar el bloque respecto de su herencia estructural.
\end{observation}

\section{Condición mínima de transición}
\begin{definition}
Una transición de nivel es admisible cuando concurren simultáneamente:
\begin{enumerate}
  \item conservación de dependencia respecto del núcleo basal integrado en Tomo I;
  \item ganancia operativa real;
  \item mantenimiento de legibilidad estructural;
  \item compatibilidad con el dominio verificable.
\end{enumerate}
\end{definition}

\begin{observation}
No toda ampliación verbal constituye transición de nivel. Debe haber mejora efectiva de tratamiento y no simple desplazamiento retórico.
\end{observation}

\section{Herencia y no refundación}
\begin{definition}
Toda transición de nivel debe poder leerse como transformación o enriquecimiento de una estructura previa, y no como refundación encubierta de primitivas ya fijadas.
\end{definition}

\begin{observation}
Cuando una pieza solo puede sostenerse rehechando lo ya definido en la base, la transición no se considera todavía disciplinada.
\end{observation}

\section{Gradiente de fortalecimiento}
\begin{definition}
Una transición de nivel puede fortalecer una formulación mediante una o más de las siguientes operaciones:
\begin{enumerate}
  \item explicitación de umbral;
  \item explicitación de criterio;
  \item explicitación de jerarquía;
  \item explicitación de compatibilidad;
  \item explicitación de condición de entrada o salida.
\end{enumerate}
\end{definition}

\begin{observation}
Estas operaciones no agotan todo fortalecimiento posible, pero bastan para fijar un primer gradiente disciplinado en el Tomo II.
\end{observation}

\section{Transición y pérdida}
\begin{definition}
Se considera fallida una transición de nivel cuando la nueva formulación pierde legibilidad, rompe herencia canónica o no aporta ganancia operativa discernible.
\end{definition}

\begin{observation}
Esta cláusula evita inflar el canon intermedio con piezas que solo aparentan mayor fuerza sin producir mejora real de organización o tratamiento.
\end{observation}

\section{Cláusula de alcance}
\begin{observation}
La transición de nivel fija una gramática de paso entre formulaciones dentro del canon intermedio. No constituye todavía una teoría completa de jerarquías ni una taxonomía final de todos los regímenes del corpus.
\end{observation}



